\subsection{Basic Facts About Modular Arithmetic}

\content{
        \begin{theorem} (Basic facts about modular arithmetic) when $n$ is a
        natural number, the following are true: 
        \begin{enumerate}
        \item $ (a+b)\mod{n} = (a\mod{n}) + (b\mod{n}) $
        \item $ (ab)\mod{n} = (a\mod{n})(b\mod{n}) $
        \item $ a^b \mod{n} = (a\mod{n})^b \mod{n} $
        \end{enumerate}
        \end{theorem}
        
        \begin{example} Compute the following:
        \begin{enumerate}
        \item $ 128\cdot 8935 \mod{5} $
        \item $ 4356^{2532} \mod{4355} $
        \item $ 4354^{2532} \mod{4355} $
        \end{enumerate}
        \end{example}
}{% presentation
\vspace{3in}
}{% nopresentation
\vspace{3in}
}{% comments
}

\content{
        \begin{example} Compute the following:
        \begin{enumerate}
        \item $ 47^{8} \mod{49} $
        \item $ 65^{5} \mod{35} $
        \end{enumerate}
        \end{example}
}{% presentation
\vspace{3in}
}{% nopresentation
\vspace{3in}
}{% comments
}

\content{
        \begin{example} Your turn! Compute the following:
        \begin{enumerate}
        \item $ 12^8 \mod{4} $
        \item $ 12^8 \mod{5} $
        \item $ 5^12 \mod{4} $
        \item $ 98^{123} \mod{99} $
        \end{enumerate}
        \end{example}
}{% presentation
}{% nopresentation
}{% comments
}

\content{
        What else can we do with modular arithmetic? 

        \begin{example} Compute each of the following:
        \begin{enumerate}
        \item $45^2\mod{20}$
        \item $(-35)^2\mod{50}$
        \end{enumerate}
        \end{example}
}{% presentation
\vspace{3in}
}{% nopresentation
\vspace{3in}
}{% comments
Notice that these give the same answer. It turns out that any computation that
we do involving 45, can instead be done with the number -35, and vice versa.
These two numbers are said to be equivalent mod 20. What other numbers are
equivalent to -35 and 45?
}

